
\documentclass{article}
\usepackage{colortbl}
\usepackage{makecell}
\usepackage{multirow}
\usepackage{supertabular}

\begin{document}

\newcounter{utterance}

\centering \large Interaction Transcript for game `cladder', experiment `full\_v1.5\_default', episode 3821 with Qwen3.6{-}27B.
\vspace{24pt}

{ \footnotesize  \setcounter{utterance}{1}
\setlength{\tabcolsep}{0pt}
\begin{supertabular}{c@{$\;$}|p{.15\linewidth}@{}p{.15\linewidth}p{.15\linewidth}p{.15\linewidth}p{.15\linewidth}p{.15\linewidth}}
    \# & \multicolumn{2}{c}{Player} && \multicolumn{2}{c}{Game Master} \\
    \hline

    \theutterance \stepcounter{utterance}  
    & & & \multicolumn{4}{p{0.6\linewidth}}{
        \cellcolor[rgb]{0.9,0.9,0.9}{
            \makecell[{{p{\linewidth}}}]{
                \texttt{\tiny{[P1$\langle$GM]}}
                \texttt{You are an expert in causal inference. The following question is not a typical commonsense query, but rather a meticulously designed question created by a professor specializing in causal inference, intended to assess the students' mastery of the course content.} \\
\texttt{Imagine a self{-}contained, hypothetical world with only the following conditions, and without any unmentioned factors or causal relationships: Gender has a direct effect on department competitiveness and peanut allergy. Department competitiveness has a direct effect on peanut allergy. For individuals who are not male, the probability of being allergic to peanuts is 48\%. For individuals who are male, the probability of being allergic to peanuts is 36\%. For individuals who are male, would it be more likely to see being allergic to peanuts if the individual had not been male?} \\
\texttt{Start your answer with 'Yes' or 'No', followed by additional reasoning or evidence to support your explanation.} \\
            }
        }
    }
     \\ \\

    \theutterance \stepcounter{utterance}  
    & \multicolumn{4}{p{0.6\linewidth}}{
        \cellcolor[rgb]{0.561,0.737,0.561}{
            \makecell[{{p{\linewidth}}}]{
                \texttt{\tiny{[P1$\rangle$GM]}}
                \texttt{Yes} \\
\\ 
\texttt{**Reasoning:**} \\
\\ 
\texttt{To answer this question, we must determine the **counterfactual probability** of peanut allergy for a male individual, had they not been male (i.e., had they been female). This requires calculating the **Average Treatment Effect (ATE)** or, more specifically, the effect of gender on peanut allergy for the treated group (males), often referred to as the **Average Treatment Effect on the Treated (ATT)**.} \\
\\ 
\texttt{Let's define the variables:} \\
\texttt{*   \$G\$: Gender (\$G=1\$ for Male, \$G=0\$ for Female)} \\
\texttt{*   \$D\$: Department Competitiveness} \\
\texttt{*   \$Y\$: Peanut Allergy (\$Y=1\$ if allergic, \$Y=0\$ if not)} \\
\\ 
\texttt{The problem states the following causal structure:} \\
\texttt{1.  \$G \textbackslash{}rightarrow D\$} \\
\texttt{2.  \$G \textbackslash{}rightarrow Y\$} \\
\texttt{3.  \$D \textbackslash{}rightarrow Y\$} \\
\\ 
\texttt{This implies that \$D\$ is a **mediator** (or confounder, depending on interpretation, but here it's a direct effect from G) in the causal path from Gender to Peanut Allergy. However, the question asks for a counterfactual: "For individuals who are male, would it be more likely to see being allergic to peanuts if the individual had not been male?"} \\
\\ 
\texttt{This is asking to compare:} \\
\texttt{*   \$P(Y=1 | G=1)\$: The observed probability of} \\
            }
        }
    }
    & & \\ \\

    \theutterance \stepcounter{utterance}  
    & & & \multicolumn{2}{p{0.3\linewidth}}{
        \cellcolor[rgb]{0.9,0.9,0.9}{
            \makecell[{{p{\linewidth}}}]{
                \texttt{\tiny{[GM$|$GM]}}
                \texttt{Yes} \\
            }
        }
    }
    & & \\ \\

    \theutterance \stepcounter{utterance}  
    & & & \multicolumn{2}{p{0.3\linewidth}}{
        \cellcolor[rgb]{0.9,0.9,0.9}{
            \makecell[{{p{\linewidth}}}]{
                \texttt{\tiny{[GM$|$GM]}}
                \texttt{yes} \\
            }
        }
    }
    & & \\ \\

    \theutterance \stepcounter{utterance}  
    & & & \multicolumn{2}{p{0.3\linewidth}}{
        \cellcolor[rgb]{0.9,0.9,0.9}{
            \makecell[{{p{\linewidth}}}]{
                \texttt{\tiny{[GM$|$GM]}}
                \texttt{game\_result = WIN} \\
            }
        }
    }
    & & \\ \\

\end{supertabular}
}

\end{document}
